\section*{Capítulo \ref{ch:b3:membrane-traffic} --- Tráfico de membranas y clasificación de proteínas}

\begin{solution}{exo:b3:membrane-traffic:1}
Tramo hidrófobo aminoterminal: \omterm{def:b3:membrane-traffic:signals}{péptido señal}, al retículo endoplasmático
(y por defecto de ahí a la superficie). PKKKRKV: \omterm{def:b3:membrane-traffic:signals}{señal de localización
nuclear}, al núcleo a través del poro. Hélice anfipática aminoterminal:
\omterm{def:b3:membrane-traffic:signals}{presecuencia} mitocondrial, a la matriz mitocondrial. KDEL: recuperación
desde el Golgi de vuelta al retículo. Manosa-6-fosfato: del Golgi trans
al \omterm{def:b3:membrane-traffic:endocytosis}{endosoma} y al \omterm{def:b3:membrane-traffic:endocytosis}{lisosoma}.
\end{solution}

\begin{solution}{exo:b3:membrane-traffic:2}
(1) Sin membranas, la cadena traducida era unos veinte residuos más larga
que la proteína secretada: se retira un segmento durante la secreción.
(2) Con membranas presentes durante la traducción, el producto tenía el
tamaño maduro y estaba protegido de la proteasa: había entrado en las
vesículas y perdido el segmento dentro de ellas. (3) Las membranas
añadidas después de la traducción no hacían ni lo uno ni lo otro: la
entrada va acoplada a la síntesis --- es cotraduccional --- y el
segmento, el \omterm{def:b3:membrane-traffic:signals}{péptido señal}, solo funciona en una cadena naciente.
\end{solution}

\begin{solution}{exo:b3:membrane-traffic:3}
Del retículo al Golgi: \omterm{def:b3:membrane-traffic:coats}{COPII}, hacia delante. Del Golgi al retículo:
\omterm{def:b3:membrane-traffic:coats}{COPI}, retrógrado. Del Golgi trans al \omterm{def:b3:membrane-traffic:endocytosis}{endosoma}: \omterm{def:b3:membrane-traffic:coats}{clatrina} con sus
adaptadores, hacia delante en dirección al \omterm{def:b3:membrane-traffic:endocytosis}{lisosoma}. De la membrana
plasmática al \omterm{def:b3:membrane-traffic:endocytosis}{endosoma}: \omterm{def:b3:membrane-traffic:coats}{clatrina}, hacia dentro (endocitosis).
\end{solution}

\begin{solution}{exo:b3:membrane-traffic:4}
Los azúcares solo se añaden dentro de la luz del retículo y del Golgi. La
cara luminal de toda vesícula sigue siendo luminal a lo largo de la
gemación y la fusión, y cuando una vesícula se fusiona con la membrana
plasmática su cara luminal pasa a ser el exterior de la célula; la
orientación no se invierte nunca, de modo que lo que se glicosiló dentro
acaba fuera.
\end{solution}

\begin{solution}{exo:b3:membrane-traffic:5}
$t^{*} = \ln(0.05/0.2)/(0.05 - 0.2) = \ln 4/0.15 = \qty{9.2}{min}$;
$B(t^{*}) = (0.2/(-0.15))(e^{-1.85} - e^{-0.46}) = 0.63$. A los
\qty{60}{min}, $A \approx 0$, $B = 1.33\,(e^{-3} - e^{-12}) = 0.066$,
luego $C = 0.93$.
\end{solution}

\begin{solution}{exo:b3:membrane-traffic:6}
El \omterm{def:b3:membrane-traffic:signals}{péptido señal} mete el extremo amino en la luz; cada tramo hidrófobo es
alternativamente una parada o un inicio de transferencia, de modo que hay
cuatro hélices transmembranarias (aproximadamente los residuos 30--52,
90--112, 150--172 y 210--232). Segmentos: extremo amino, luminal; 52--90,
citosólico; 112--150, luminal; 172--210, citosólico; extremo carboxilo,
luminal. Glicosilación solo en los segmentos luminales --- el extremo
amino, 112--150 y el extremo carboxilo ---, que pasan a ser
extracelulares en la superficie.
\end{solution}

\begin{solution}{exo:b3:membrane-traffic:7}
$J_{\max} = \num{30000}\times 0.25\times 0.05/0.30 = 1250$ por minuto;
fracción de superficie $k_{r}/(k_{i} + k_{r}) = 0.05/0.30 = 0.17$.
Duplicar $R$ duplica $J$ exactamente; el cuello de botella es la
velocidad de regreso (la mayoría de los receptores están dentro), y
duplicar $k_{r}$ daría $\num{30000}\times 0.25
\times 0.1/0.35 \approx 2140$, un factor $1.7$.
\end{solution}

\begin{solution}{exo:b3:membrane-traffic:8}
(a) Sin receptor: las hidrolasas marcadas siguen la ruta por defecto y se
secretan; los \omterm{def:b3:membrane-traffic:endocytosis}{lisosomas} se quedan sin enzimas y se llenan de material sin
digerir --- una enfermedad de depósito. (b) Sin fosfotransferasa: el
mismo fenotipo por otra lesión (la enfermedad de las células~I), con las
enzimas sin marcar y secretadas. (c) \omterm{def:b3:membrane-traffic:endocytosis}{Endosomas} neutralizados: el receptor
no puede soltar su carga, de modo que no se recicla y las hidrolasas se
entregan mal o se secretan; los receptores de LDL y otros también dejan
de ciclar.
\end{solution}

\begin{solution}{exo:b3:membrane-traffic:9}
La tirosina pertenece al \omterm{def:b3:bioinformatics:motif}{motivo} de la cola que reconoce el adaptador de
\omterm{def:b3:membrane-traffic:coats}{clatrina}; sin ella el receptor no se reúne en las fositas recubiertas.
Los receptores unen la LDL con normalidad pero están repartidos
uniformemente por la superficie, excluidos de las fositas, y el ligando
se queda fuera --- una enfermedad de clasificación, no de unión.
\end{solution}

\begin{solution}{exo:b3:membrane-traffic:10}
Con $k_{1} = k_{2} = k$, pruébese $B = kt\,e^{-kt}$: $\mathrm{d}B/\mathrm{d}t
= k e^{-kt} - k^{2}t e^{-kt} = kA - kB$, como se pedía, con $B(0) = 0$.
El máximo está donde $1 - kt = 0$: $t^{*} = 1/k$, $B = e^{-1} = 0.37$.
En general
$C(t) = 1 - (k_{2}e^{-k_{1}t} - k_{1}e^{-k_{2}t})/(k_{2} - k_{1})$,
que es simétrico al intercambiar $k_{1}$ y $k_{2}$: la curva de
los gránulos por sí sola no puede decir cuál de los dos compartimentos es
el rápido.
\end{solution}

\begin{solution}{exo:b3:membrane-traffic:11}
Cada vesícula tiene un área $4\pi(50\,\text{nm})^{2} = \qty{0.031}{\micro
m^2}$; mil por minuto son \qty{31}{\micro m^2/min}, de modo que toda la
superficie en $3000/31 \approx \qty{95}{min}$. La célula no se encoge
porque la membrana vuelve por exocitosis (vesículas de reciclado) al
mismo ritmo --- un equilibrio de flujos. Cada vesícula contiene
$5\times 10^{-19}$ L; mil por minuto durante una hora son
$3\times 10^{-14}$ L, alrededor del \qty{1.5}{\%} del volumen de una
célula de \qty{2}{pL} por hora.
\end{solution}

\begin{solution}{exo:b3:membrane-traffic:12}
La toxina botulínica actúa en el terminal motor: no se libera
acetilcolina, el músculo no recibe orden y queda fláccido --- parálisis
flácida. La toxina tetánica sube por el axón motor y pasa a las
interneuronas inhibidoras que normalmente frenan a las motoneuronas; con
su \omterm{def:b3:membrane-traffic:fusion}{SNARE} cortada no se libera glicina ni GABA, las motoneuronas disparan
sin oposición y todos los músculos se contraen --- parálisis espástica.
Unos pocos nanogramos de toxina botulínica inyectados en un músculo
hiperactivo lo silencian localmente durante meses (hasta que se fabrica
\omterm{def:b3:membrane-traffic:fusion}{SNARE} nueva y el terminal se recupera): un tratamiento para distonías,
espasticidad, estrabismo y arrugas.
\end{solution}

\begin{solution}{pb:b3:membrane-traffic:1}
\textbf{1.} $10^{7}\times \qty{25000}{Da}\times 1.66\times 10^{-24}$ g
$= \qty{0.42}{pg}$ por minuto, \qty{600}{pg} al día --- el doble de su
propio contenido proteico.
\textbf{2.} Volumen de la vesícula,
$\tfrac{4}{3}\pi\,35^{3} = 1.8\times 10^{5}$ nm$^{3}$, y el
\qty{30}{\%} de él, $5.4\times 10^{4}$; una
molécula de enzima, $\tfrac{4}{3}\pi\,2^{3} = 34$ nm$^{3}$: unas
\num{1600} moléculas por vesícula; $10^{7}/1600 \approx 6200$ vesículas
por minuto.
\textbf{3.} Área $4\pi\,35^{2} = \qty{0.0154}{\micro m^2}$ cada una,
\qty{95}{\micro
m^2} por minuto: los \qty{6000}{\micro m^2} del retículo en alrededor de
una hora si no volviera nada.
\textbf{4.} $\qty{95}{\micro m^2} = 9.5\times 10^{7}$ nm$^{2}$, con dos
monocapas a \qty{0.7}{nm^2}: $2.7\times 10^{8}$ lípidos por minuto. La
célula devuelve membrana con vesículas \omterm{def:b3:membrane-traffic:coats}{COPI} y reciclado, y sintetiza
lípido para reponer lo que se va para siempre en los gránulos.
\textbf{5.} Volumen del gránulo, $5.2\times 10^{8}$ nm$^{3}$, y el
\qty{30}{\%} de él, $1.6\times 10^{8}$: $4.7\times 10^{6}$ moléculas por
gránulo; $6\times
10^{8}$ moléculas por hora llenan unos $130$ gránulos.
\textbf{6.} Cada gránulo añade $\pi d^{2} = \qty{3.1}{\micro m^2}$; $300$
añaden \qty{940}{\micro m^2} a \qty{50}{\micro m^2}: veinte veces más. La
membrana ha de recuperarse por endocitosis compensadora en cuestión de
minutos.
\textbf{7.} $t^{*} = \ln 2/0.05 = \qty{13.9}{min}$, $B = 0.50$.
\textbf{8.} $A = e^{-3} = 0.05$; $B = 2(e^{-1.5} - e^{-3}) = 0.35$;
$C = 0.60$.
\textbf{9.} \qty{10}{min} y \qty{20}{min}; total medio \qty{30}{min}.
\textbf{10.} $t^{*} = \ln(0.1)/(-0.09) = \qty{25.6}{min}$; $B(t^{*}) =
(0.1/(-0.09))(e^{-2.56} - e^{-0.256}) = 0.77$: el Golgi se hincha,
atiborrado con tres cuartas partes de la carga.
\textbf{11.} En el máximo, $k_{1}e^{-k_{1}t^{*}} = k_{2}e^{-k_{2}t^{*}}$;
al sustituir en $B$ resulta $B_{\max} = (k_{1}/k_{2})^{\,k_{2}/(k_{2} -
k_{1})}$: la potencia es $k_{2}/(k_{2} - k_{1})$. Para $k_{1} > k_{2}$ el
exponente es negativo y la base mayor que $1$; para $k_{1} < k_{2}$ la
base es menor que $1$ y el exponente positivo --- en ambos casos el valor
es menor que $1$ (compruébese: aquí $2^{-1} = 0.5$).
\textbf{12.} Para $k_{1} \gg k_{2}$, $B(t) \to e^{-k_{2}t}$: la marca
está en el Golgi de inmediato y sale a $k_{2}$ --- el paso del retículo
es invisible y la curva del Golgi es una simple caída.
\textbf{13.} $\theta = 2/(2+2) = 0.5$; $J = \num{50000}\times 0.2\times
0.1\times 0.5/(0.1 + 0.1) = 2500$ partículas por minuto, \num{150000} por
hora.
\textbf{14.} $1.5\times 10^{5}\times 1500 = 2.3\times 10^{8}$ moléculas
de colesterol por hora; $5\times 10^{9}/2.3\times 10^{8} \approx 22$
horas.
\textbf{15.} Fracción en superficie, $k_{r}/(k_{r} + k_{i}\theta) = 0.5$;
un viaje de ida y vuelta cuesta
$1/(k_{i}\theta) + 1/k_{r} = 10 + 10 = \qty{20}{min}$: unos $60$ viajes
en veinte horas.
\textbf{16.} $J = 1250$ por minuto. Por debajo de la saturación
$J \propto R\,[L]$, de modo que con la mitad de los receptores la LDL
plasmática ha de duplicarse para el mismo aclaramiento.
\textbf{17.} $R$ vuelve a \num{50000}: se restaura el aclaramiento y la
LDL plasmática baja hacia lo normal --- el efecto de la estatina.
\textbf{18.} Sin receptores, la LDL solo se aclara por vías inespecíficas
y lentas, de modo que circula durante días en lugar de horas, se oxida y
la captan los receptores basureros de los macrófagos, que se convierten
en las células espumosas de las placas ateroscleróticas; y las estatinas,
que actúan induciendo receptores, no tienen nada que inducir.
\textbf{19.} $V = \tfrac{4}{3}\pi(0.25)^{3} = \qty{0.065}{\micro m^3} =
6.5\times 10^{-17}$ L. A pH $7.2$: $6.3\times 10^{-8}\times 6.5\times
10^{-17} = 4\times 10^{-24}$ mol --- unos $2$ protones libres. A pH
$4.7$: $2\times 10^{-5}\times 6.5\times 10^{-17} = 1.3\times 10^{-21}$
mol, unos $780$.
\textbf{20.} $50\times 780 \approx \num{39000}$ protones; $50$ bombas a
$200$ por segundo mueven \num{10000} por segundo: unos \qty{4}{s}.
\textbf{21.} Meter carga positiva deja la luz positiva, y unos pocos
miles de cargas detendrían la bomba; entra cloruro por un canal (o salen
cationes) para neutralizar la carga.
\textbf{22.} Las enzimas inactivas a pH neutro no hacen daño si un
\omterm{def:b3:membrane-traffic:endocytosis}{lisosoma} tiene una fuga o si se clasifican mal hacia el citosol o la
superficie: la exigencia de acidez confina la digestión al compartimento
construido para ella.
\textbf{23.} Ambos han de llevar un sensor de pH --- histidinas, cuya
protonación cerca de pH~6 cambia la superficie de unión: el receptor de
LDL pliega un dominio sobre su propio sitio de unión al protonarse, y el
receptor de manosa-6-fosfato pierde afinidad de forma análoga.
\textbf{24.} $10^{7.2 - 4.7} = 10^{2.5} \approx 300$ veces más
concentrada; la base atrapada consume protones, sube el pH lisosómico e
inactiva las hidrolasas --- que es como la cloroquina mata a los
parásitos del paludismo en su vacuola alimentaria y como bloquea la
\omterm{def:b3:membrane-traffic:endocytosis}{autofagia} en los experimentos.
\textbf{25.} Máximo del Golgi a los \qty{14}{min}; unas \num{6200}
vesículas salen del retículo por minuto; y \num{150000} partículas de LDL
captadas por hora.
\end{solution}
